% =========================================================
\subsubsection{Multiple Metrics on the Same Set}
\label{sec:multiple-metrics}
% =========================================================

A metric is not a property of the underlying set --- it is an additional
structure imposed on it.
The same set $X$ can carry many different metrics simultaneously, each
producing a different geometry, different unit balls, and potentially
different notions of convergence.
This section develops three related ideas:
(1) the familiar plane $\mathbb{R}^2$ with three standard metrics;
(2) the product metric, which builds a single metric on $X \times Y$ from
separate metrics on $X$ and $Y$;
(3) the interpretation of a product metric as assigning independent
coordinate metrics to each axis.

% ---- Toolkit ----
\begin{tcolorbox}[colback=gray!6, colframe=gray!40, arc=2pt,
  left=6pt, right=6pt, top=4pt, bottom=4pt,
  title={\small\textbf{Multiple Metrics --- Quick Reference}},
  fonttitle=\small\bfseries]
\small
\begin{tabular}{l l l l}
\toprule
\textbf{Label} & \textbf{Object} & \textbf{Key fact} & \textbf{Detail} \\
\midrule
D\ref{def:r2-three-metrics}
  & $(\mathbb{R}^2, d_1)$, $(\mathbb{R}^2, d_2)$, $(\mathbb{R}^2, d_\infty)$
  & All equivalent; different unit balls
  & \hyperref[def:r2-three-metrics]{$\downarrow$} \\
D\ref{def:product-metric}
  & Product metric $(X \times Y,\, d)$
  & Three standard variants; all equivalent
  & \hyperref[def:product-metric]{$\downarrow$} \\
D\ref{def:product-coord-metrics}
  & Coordinate metric spaces
  & Independent metrics per axis
  & \hyperref[def:product-coord-metrics]{$\downarrow$} \\
P\ref{prop:product-convergence}
  & Convergence in product metric
  & $(x_n,y_n)\to(x,y)$ iff $x_n\to x$ and $y_n\to y$
  & \hyperref[prop:product-convergence]{$\downarrow$} \\
\bottomrule
\end{tabular}
\end{tcolorbox}

% =========================================================
% 1.  R^2 with d1, d2, d_inf
% =========================================================

\begin{tcolorbox}[colback=propbox, colframe=propborder, arc=2pt,
  left=6pt, right=6pt, top=4pt, bottom=4pt,
  title={\small\textbf{Definition ($\mathbb{R}^2$ with Three Standard Metrics)}},
  fonttitle=\small\bfseries]
\begin{definition}[$\mathbb{R}^2$ with Three Standard Metrics]
\label{def:r2-three-metrics}
Let $p = (x_1, x_2)$ and $q = (y_1, y_2)$ be points in $\mathbb{R}^2$.
The three standard metrics on $\mathbb{R}^2$ are:

\medskip
\noindent\textbf{Taxicab metric} $d_1$:
\[
d_1(p,\,q) \;:=\; |x_1 - y_1| + |x_2 - y_2|.
\]

\noindent\textbf{Euclidean metric} $d_2$:
\[
d_2(p,\,q) \;:=\; \sqrt{(x_1-y_1)^2 + (x_2-y_2)^2}.
\]

\noindent\textbf{Maximum metric} $d_\infty$:
\[
d_\infty(p,\,q) \;:=\; \max\bigl(|x_1 - y_1|,\; |x_2 - y_2|\bigr).
\]

\medskip
\noindent\textbf{Convergence.}\quad
A sequence $(p^{(n)}) = (x_1^{(n)}, x_2^{(n)})$ in $\mathbb{R}^2$
converges to $p = (x_1, x_2)$ under any of these three metrics if and
only if it converges coordinatewise:
\[
\forall\,\varepsilon > 0,\;\exists N \in \mathbb{N},\;
\forall n \ge N:\quad
\bigl|x_1^{(n)} - x_1\bigr| < \varepsilon
\;\;\text{and}\;\;
\bigl|x_2^{(n)} - x_2\bigr| < \varepsilon.
\]
\end{definition}
\end{tcolorbox}

\begin{remark*}[Fully quantified form]
\begin{align*}
\forall\, p,q \in \mathbb{R}^2: \quad
d_1(p,q) &\;=\; |x_1-y_1| + |x_2-y_2|, \\
d_2(p,q) &\;=\; \bigl((x_1-y_1)^2 + (x_2-y_2)^2\bigr)^{1/2}, \\
d_\infty(p,q) &\;=\; \max\bigl(|x_1-y_1|,\,|x_2-y_2|\bigr).
\end{align*}
Each satisfies all four metric axioms; the non-trivial one is the triangle
inequality, which for $d_1$ follows from $|a+b| \le |a|+|b|$ applied
coordinatewise, and for $d_2$ is Minkowski's inequality.
\end{remark*}

\begin{remark*}[Metric inequalities]
For all $p, q \in \mathbb{R}^2$, the three metrics satisfy
\[
d_\infty(p,q)
\;\le\;
d_2(p,q)
\;\le\;
d_1(p,q)
\;\le\;
2\,d_\infty(p,q).
\]
This chain of inequalities shows the metrics are \emph{strongly equivalent}
(Definition~\ref{def:equivalent-metrics}): they produce the same open sets,
the same convergent sequences, and the same compact sets.
The constant $2$ in the last inequality is sharp and depends on the
dimension; in $\mathbb{R}^n$ it becomes $\sqrt{n}$ for the middle pair
and $n$ for the outer pair.
\end{remark*}

\begin{remark*}[Unit ball geometry]
Although the metrics are equivalent, their unit balls
$B_1(0) = \{q \in \mathbb{R}^2 : d(0, q) < 1\}$ look very different:
$d_1$ gives a diamond (rotated square),
$d_2$ gives a disc,
$d_\infty$ gives a square aligned with the axes.
See Figure~\ref{fig:unit-balls} and Figure~\ref{fig:unit-balls-comparison}.
\end{remark*}

% =========================================================
% 2.  Product metric
% =========================================================

\begin{tcolorbox}[colback=propbox, colframe=propborder, arc=2pt,
  left=6pt, right=6pt, top=4pt, bottom=4pt,
  title={\small\textbf{Definition (Product Metric Space)}},
  fonttitle=\small\bfseries]
\begin{definition}[Product Metric]
\label{def:product-metric}
Let $(X, d_X)$ and $(Y, d_Y)$ be metric spaces.
The \emph{product metric space} is the Cartesian product $X \times Y$
equipped with one of the following equivalent metrics.
For points $(x_1, y_1),\, (x_2, y_2) \in X \times Y$:

\medskip
\noindent\textbf{Sum metric} ($\ell^1$ product):
\[
d_{\mathrm{sum}}\bigl((x_1,y_1),(x_2,y_2)\bigr)
\;:=\; d_X(x_1,x_2) \;+\; d_Y(y_1,y_2).
\]

\noindent\textbf{Max metric} ($\ell^\infty$ product):
\[
d_{\max}\bigl((x_1,y_1),(x_2,y_2)\bigr)
\;:=\; \max\bigl(d_X(x_1,x_2),\; d_Y(y_1,y_2)\bigr).
\]

\noindent\textbf{Euclidean product metric} ($\ell^2$ product):
\[
d_{\mathrm{euc}}\bigl((x_1,y_1),(x_2,y_2)\bigr)
\;:=\; \sqrt{d_X(x_1,x_2)^2 \;+\; d_Y(y_1,y_2)^2}.
\]

\medskip
\noindent\textbf{Convergence.}\quad
A sequence $\bigl((x_n, y_n)\bigr)_{n \ge 1}$ in $X \times Y$ converges
to $(x, y)$ under any product metric if and only if:
\[
\forall\,\varepsilon > 0,\;\exists N \in \mathbb{N},\;
\forall n \ge N:\quad
d_X(x_n, x) < \varepsilon
\;\;\text{and}\;\;
d_Y(y_n, y) < \varepsilon.
\]
The two coordinates converge \emph{independently} in their own metrics.
\end{definition}
\end{tcolorbox}

\begin{remark*}[Fully quantified form --- sum metric]
\[
\forall\,(x_1,y_1),\,(x_2,y_2) \in X \times Y, \quad
d_{\mathrm{sum}}\bigl((x_1,y_1),(x_2,y_2)\bigr)
\;=\; d_X(x_1,x_2) + d_Y(y_1,y_2) \;\in\; [0,\infty).
\]
The four axioms follow directly from those of $d_X$ and $d_Y$; for the
triangle inequality:
\begin{align*}
d_{\mathrm{sum}}\bigl((x_1,y_1),(x_3,y_3)\bigr)
&= d_X(x_1,x_3) + d_Y(y_1,y_3) \\
&\le \bigl(d_X(x_1,x_2) + d_X(x_2,x_3)\bigr)
   + \bigl(d_Y(y_1,y_2) + d_Y(y_2,y_3)\bigr) \\
&= d_{\mathrm{sum}}\bigl((x_1,y_1),(x_2,y_2)\bigr)
 + d_{\mathrm{sum}}\bigl((x_2,y_2),(x_3,y_3)\bigr).
\end{align*}
\end{remark*}

\begin{remark*}[All three product metrics are equivalent]
For all $(x_1,y_1), (x_2,y_2) \in X \times Y$,
\[
d_{\max} \;\le\; d_{\mathrm{euc}} \;\le\; d_{\mathrm{sum}}
\;\le\; 2\,d_{\max}.
\]
They therefore induce the same topology, the same convergent sequences,
and the same Cauchy sequences on $X \times Y$.
The choice is a matter of convenience; $d_{\max}$ is usually the simplest
for proofs, $d_{\mathrm{euc}}$ is natural when $X = Y = \mathbb{R}$.
\end{remark*}

\begin{remark*}[Completeness is inherited]
If $(X, d_X)$ and $(Y, d_Y)$ are both complete then $(X \times Y, d)$
is complete under any product metric.
This is because a Cauchy sequence in $X \times Y$ has Cauchy coordinate
projections, which converge by completeness of $X$ and $Y$ separately.
\end{remark*}

% =========================================================
% 3.  Coordinate metric spaces
% =========================================================

\begin{tcolorbox}[colback=propbox, colframe=propborder, arc=2pt,
  left=6pt, right=6pt, top=4pt, bottom=4pt,
  title={\small\textbf{Definition (Coordinate Metric Spaces)}},
  fonttitle=\small\bfseries]
\begin{definition}[Coordinate Metric Spaces]
\label{def:product-coord-metrics}
Let $(X, d_X)$ and $(Y, d_Y)$ be metric spaces, and let
$(X \times Y, d)$ be their product (Definition~\ref{def:product-metric}).
We say that:
\begin{itemize}
  \item the \emph{$X$-coordinate} of a point $(x, y) \in X \times Y$ lives
        in the metric space $(X, d_X)$;
  \item the \emph{$Y$-coordinate} lives in $(Y, d_Y)$;
  \item $(X, d_X)$ and $(Y, d_Y)$ are the \emph{coordinate metric spaces}
        of the product.
\end{itemize}
When $X = Y = \mathbb{R}$, this means the $x$-axis and $y$-axis of
$\mathbb{R}^2$ carry independent metrics $d_X$ and $d_Y$ which need not
be the same.

\medskip
\noindent\textbf{Canonical example.}\quad
Take $d_X(x_1,x_2) := |x_1 - x_2|$ (standard metric on $\mathbb{R}$) and
$d_Y(y_1,y_2) := \mathbf{1}_{y_1 \neq y_2}$ (discrete metric on $\mathbb{R}$).
The product space $(\mathbb{R} \times \mathbb{R},\, d_{\mathrm{sum}})$ has
metric:
\[
d\bigl((x_1,y_1),(x_2,y_2)\bigr)
\;:=\; |x_1 - x_2| \;+\; \mathbf{1}_{y_1 \neq y_2}.
\]
Horizontal distances are continuous; any nonzero vertical displacement costs
exactly $1$, regardless of its size.

\medskip
\noindent\textbf{Convergence.}\quad
$\bigl((x_n, y_n)\bigr) \to (x, y)$ in this space if and only if:
\[
|x_n - x| \to 0 \;\text{ in } \mathbb{R}
\quad\text{and}\quad
y_n = y \;\text{ for all sufficiently large } n.
\]
\end{definition}
\end{tcolorbox}

\begin{remark*}[Fully quantified form --- standard $\times$ discrete]
\[
\forall\,(x_1,y_1),(x_2,y_2) \in \mathbb{R}^2, \quad
d\bigl((x_1,y_1),(x_2,y_2)\bigr)
\;=\; |x_1 - x_2| + \mathbf{1}_{y_1 \neq y_2} \;\in\; [0,\infty).
\]
Positivity: $d = 0$ iff $|x_1-x_2| = 0$ and $\mathbf{1}_{y_1\neq y_2} = 0$
iff $x_1 = x_2$ and $y_1 = y_2$. Symmetry and triangle inequality hold
because both $|\cdot|$ and $\mathbf{1}_{\neq}$ individually satisfy them.
\end{remark*}

\begin{remark*}[Unit ball of the standard $\times$ discrete product]
The unit ball around the origin is
\[
B_1\bigl((0,0)\bigr)
= \bigl\{\,(x,y) \in \mathbb{R}^2 :
|x| + \mathbf{1}_{y \neq 0} < 1\,\bigr\}.
\]
If $y = 0$: the condition is $|x| < 1$, giving the open interval $(-1,1)$
on the $x$-axis.
If $y \neq 0$: the condition is $|x| + 1 < 1$, i.e.\ $|x| < 0$, which is
impossible.
The unit ball is therefore exactly the open segment
$\{(x,0) : |x| < 1\}$ --- a one-dimensional interval, not a
two-dimensional region.
See Figure~\ref{fig:unit-balls-comparison}.
\end{remark*}

\begin{remark*}[Why coordinate metrics need not agree]
In the Euclidean plane, the $x$-axis and $y$-axis both carry the standard
metric $|{\cdot}|$ and distances in both directions are interchangeable.
The product construction shows this is a special case: there is no
requirement that the two coordinate metrics be the same or even of the
same type.
Choosing different metrics per coordinate produces spaces with genuinely
anisotropic geometry --- distances in one direction are measured on a
completely different scale or topology than distances in another.
\end{remark*}

\begin{remark*}[Common error]
It is tempting to say ``the $x$-axis has metric $d_X$'' as if the metric
were a property of the axis itself.
It is not: the metric is a property of the entire space $X \times Y$.
The product construction \emph{assigns} a role to each coordinate, but
the distance function $d$ is defined on the whole of
$(X \times Y) \times (X \times Y)$, not separately on $X$ and $Y$.
\end{remark*}

% =========================================================
% 4.  Proposition: convergence in product
% =========================================================

\begin{proposition}[\texorpdfstring{%
  \hyperref[prf:product-convergence]{Convergence in the Product Metric}%
}{Convergence in the Product Metric}]
\label{prop:product-convergence}
Let $(X, d_X)$ and $(Y, d_Y)$ be metric spaces with product metric $d$
(any of the three variants).
A sequence $\bigl((x_n, y_n)\bigr)_{n \ge 1}$ in $X \times Y$ satisfies
\[
d\bigl((x_n,y_n),(x,y)\bigr) \to 0
\;\iff\;
d_X(x_n,x) \to 0 \;\text{ and }\; d_Y(y_n,y) \to 0.
\]
\end{proposition}

\begin{remark*}[Proof shape]
For the sum metric: $d_{\mathrm{sum}} = d_X + d_Y \ge d_X$ and $\ge d_Y$,
so $d_{\mathrm{sum}} \to 0$ forces both $d_X \to 0$ and $d_Y \to 0$.
Conversely, if both tend to $0$ then so does their sum.
For $d_{\max}$: identical argument with $\max$ in place of sum.
For $d_{\mathrm{euc}}$: use $d_X, d_Y \le d_{\mathrm{euc}} \le d_X + d_Y$.
\end{remark*}

\begin{remark*}[Consequence]
Convergence in the product decouples into independent convergence in each
coordinate. This means all three product metrics produce identical convergent
sequences, confirming their equivalence without requiring a direct norm comparison.
\end{remark*}